\section{Thor--FPGA System Integration}
\label{sec:system}

\subsection{Prototype and Feature Transport}
Figure~\ref{fig:prototype} presents the physical platform. The prototype
combines Jetson AGX Thor and a Stratix 10 MX development
board~\cite{nvidiathorspec,intels10mxboard,intels10mxproduct}. The current
decoder endpoint uses one 10GbE lane. Other physical lanes do not imply
aggregate decoder bandwidth unless they are actually included in the
evaluated data path.

Thor packs the boundary tensor into raw Ethernet frames. The custom header
includes a job identifier and chunk identifier, and each data chunk carries
8,192 payload bytes. The current interface uses 222 chunks for the input
tensor and 111 chunks for the returned depth map. Weights are uploaded before
inference and retained for reuse. Cold-start loading and steady-state timing
will be reported separately.

The receive bridge checks frame metadata, records received chunks, and writes
payloads into HBM. Once an input is complete, it starts the decoder. After
decoder completion, the transmit path reads the output tensor and returns its
chunks to Thor. This coarse model boundary limits the interface to features,
results, and job control; its cost remains part of the system measurement,
as required when assessing split inference~\cite{matsubara2019split}.

\subsection{Cross-Frame Execution Schedule}
\label{sec:overlap}
Let $E$ denote Thor encoding and $O$ the complete FPGA offload path, including
packing, transfer, decoding, and result handling. The intended schedule
overlaps $E_{i+1}$ with $O_i$ (Fig.~\ref{fig:timeline}). A conservative
implementation needs only one active FPGA job: host-side buffers retain the
next encoded tensor until the accelerator can accept it. It does not require
simultaneous execution of two decoder jobs.

\begin{figure}[h]
  \centering
  \begin{tikzpicture}[x=.78cm,y=.6cm,font=\scriptsize,>=Latex]
  \node[anchor=west,font=\footnotesize\bfseries] at (0,0) {Serial split};
  \node[anchor=east] at (1.6,-.8) {Thor};
  \node[anchor=east] at (1.6,-1.6) {Offload};
  \draw[fill=paperblue!17,draw=paperblue] (1.8,-1.1) rectangle (3.8,-.5);
  \node at (2.8,-.8) {$E_i$};
  \draw[fill=papergreen!17,draw=papergreen] (3.8,-1.9) rectangle (6.2,-1.3);
  \node at (5,-1.6) {$O_i$};
  \draw[fill=paperblue!17,draw=paperblue] (6.2,-1.1) rectangle (8.2,-.5);
  \node at (7.2,-.8) {$E_{i+1}$};
  \node[anchor=west,font=\footnotesize\bfseries] at (0,-2.65) {Cross-frame overlap (proposed)};
  \node[anchor=east] at (1.6,-3.5) {Thor};
  \node[anchor=east] at (1.6,-4.3) {Offload};
  \foreach \x/\lab in {1.8/i,3.8/{i+1},6.2/{i+2}} {
    \draw[fill=paperblue!17,draw=paperblue] (\x,-3.8) rectangle ++(2,.6);
    \node at (\x+1,-3.5) {$E_{\lab}$};}
  \draw[fill=papergreen!17,draw=papergreen] (3.8,-4.6) rectangle (6.2,-4);
  \node at (5,-4.3) {$O_i$};
  \draw[fill=papergreen!17,draw=papergreen] (6.2,-4.6) rectangle (8.6,-4);
  \node at (7.4,-4.3) {$O_{i+1}$};
  \draw[->,draw=black!50] (1.8,-5.1)--(8.6,-5.1) node[right]{time};
\end{tikzpicture}

  \caption{Proposed cross-frame overlap. Encoding the next frame proceeds
  while the previous frame traverses the offload path. Block lengths are
  schematic, not measured. One active FPGA job is sufficient when the next
  feature tensor waits in a separate host buffer.}
  \Description{Serial and overlapped schedules, with encoding and the complete
  offload service on separate rows.}
  \label{fig:timeline}
\end{figure}

With independent service stages and adequate buffering, the ideal period is
\begin{equation}
  T_{\mathrm{period}}\approx\max(T_E,T_O),\qquad
  L_{\mathrm{frame}}\approx T_E+T_O.
  \label{eq:pipeline}
\end{equation}
Queueing and shared-resource contention can increase both quantities. The
30-FPS target requires an observed period of at most 33.3\,ms; it does not
require the same bound on single-frame latency. Similar use of concurrent
processing stages appears in accelerator scheduling
~\cite{symons2022stream,li2025gaurast}; here the schedule spans two devices.

\subsection{Buffer Ownership and Completion}
The runtime must associate every returned tensor with its input job and
retain each buffer until its consumer finishes. Missing chunks, incomplete
outputs, or timeouts cannot count as completed frames. Host preparation and
result assembly may contend with encoding, so the final timeline will use
measured timestamps rather than assume perfect overlap. Cross-frame
scheduling is specified here as the system design; the completed-frame
experiment remains pending.
